\documentclass[12pt,a4paper]{article}
\usepackage{fontspec}
\IfFontExistsTF{Times New Roman}{\setmainfont{Times New Roman}}{\setmainfont{TeX Gyre Termes}}
\IfFontExistsTF{Arial}{\setsansfont{Arial}}{\setsansfont{TeX Gyre Heros}}
\usepackage[left=2.7cm,right=2.2cm,top=2.4cm,bottom=2.4cm]{geometry}
\usepackage{amsmath,amssymb,graphicx,booktabs,array,longtable,tabularx,float,setspace}
\usepackage[unicode,hidelinks]{hyperref}
\usepackage{caption}
\usepackage{indentfirst}
\captionsetup{font=small,labelfont=bf}
\renewcommand{\figurename}{Hình}
\renewcommand{\tablename}{Bảng}
\renewcommand{\contentsname}{MỤC LỤC}
\renewcommand{\refname}{TÀI LIỆU THAM KHẢO}
\setlength{\parindent}{1cm}
\setlength{\parskip}{4pt}
\setstretch{1.15}
\emergencystretch=2em
\newcommand{\fig}[3]{\begin{figure}[H]\centering\includegraphics[width=#2\linewidth]{images/#1}\caption{#3}\end{figure}}
\begin{document}
\begin{titlepage}
\centering
{\large\bfseries ĐẠI HỌC PHENIKAA\\TRƯỜNG CÔNG NGHỆ THÔNG TIN PHENIKAA\par}
\vspace{.7cm}
\includegraphics[width=3.2cm]{images/phenikaa.png}
\vspace{.65cm}

{\LARGE\bfseries BÁO CÁO BÀI TẬP LỚN\par}
\vspace{.4cm}
{\large\bfseries MÔN: MẠNG NƠRON VÀ HỌC SÂU\par}
\vspace{.65cm}
{\Large\bfseries AdaZero-DCE: Điều chỉnh số vòng lặp\\đường cong thích ứng cho tăng cường\\ảnh thiếu sáng\par}
\vspace{.9cm}
\begin{tabular}{ll}
\textbf{Giảng viên:} & TS. Đặng Thị Thúy An\\[5pt]
\textbf{Sinh viên:} & Trần Minh Sang -- 24100012\\[5pt]
& Ngô Quang Thiện -- 24102651\\[5pt]
\textbf{Lớp học phần:} & Mạng nơron và học sâu-1-1-26(N02)
\end{tabular}
\vspace{.7cm}

\small Mã nguồn: \url{https://github.com/ZesyC/AdaZero-DCE}
\vfill
{\large Hà Nội, tháng 10 năm 2026}
\end{titlepage}
\pagenumbering{roman}
\section*{TÓM TẮT}
Tăng cường ảnh thiếu sáng bằng Zero-DCE thường sử dụng cố định tám bước biến đổi đường cong. Tuy nhiên, số bước này có thể làm ảnh tăng sáng vượt mức cần thiết trong một số điều kiện. Báo cáo trình bày AdaZero-DCE, một phương pháp dự đoán số vòng lặp từ đặc trưng độ sáng trong khi giữ nguyên toàn bộ trọng số DCE-Net. Bộ hồi quy MLP có cấu trúc 10--128--64--1, gồm 9729 tham số, sử dụng các thống kê độ sáng và histogram của ảnh đầu vào. Cơ chế bảo vệ giữ tám vòng cho ảnh có độ sáng trung bình nhỏ hơn 0.10.

Bộ hồi quy được huấn luyện trên 1940 mẫu tạo từ 485 cặp ảnh LOLv1, gồm ảnh thiếu sáng gốc và ba mức giả lập từ ảnh tham chiếu. Thực nghiệm sử dụng ba seed và so sánh số vòng thích ứng với fixed-8 bằng cùng cách tính PSNR và SSIM. Trên eval15, tại hai mức giả lập GT nhân 0.45 và 0.6, PSNR trung bình tăng lần lượt 1.08 và 3.26 dB; SSIM cũng tăng. Trên ảnh gốc, kết quả bằng baseline. Đánh giá khái quát trên 100 cặp LOLv2-real, không huấn luyện lại, cho thấy PSNR giảm khoảng 0.20 dB và SSIM tăng nhẹ khoảng 0.001555.

Kết quả cho thấy cách chọn số vòng nhẹ có ích trên một số điều kiện giả lập, nhưng chưa chứng minh được cải thiện trên ảnh vừa tối chụp thật hoặc khả năng khái quát tốt ngoài LOLv1. Đóng góp của báo cáo là cách triển khai gọn và phép so sánh định lượng có kiểm soát; ý tưởng điều chỉnh thích ứng đã xuất hiện trong các nghiên cứu trước.

\textbf{Từ khóa:} tăng cường ảnh thiếu sáng; Zero-DCE; số vòng lặp thích ứng; hồi quy độ sáng; LOLv1; LOLv2-real.
\clearpage
\begingroup\small\setstretch{0.93}
\tableofcontents
\endgroup
\clearpage
\pagenumbering{arabic}
\section{Giới thiệu bài toán}
\subsection{Đặt vấn đề}
Ảnh thu nhận trong điều kiện thiếu sáng thường có độ tương phản thấp, vùng tối khó quan sát và màu sắc không rõ ràng. Tăng cường ảnh thiếu sáng nhằm tạo ảnh dễ quan sát hơn nhưng vẫn giữ cấu trúc và nội dung của cảnh. Việc tăng độ sáng đồng loạt có thể làm các vùng vốn đủ sáng trở nên quá sáng, trong khi chi tiết tại vùng tối chưa chắc được khôi phục. Vì vậy, mức độ điều chỉnh cần phù hợp với từng ảnh đầu vào.

Zero-DCE \cite{zerodce} tiếp cận bài toán bằng cách ước lượng các đường cong biến đổi điểm ảnh. Mạng DCE-Net sinh bản đồ hệ số, sau đó áp dụng nhiều bước đường cong để tạo ảnh tăng cường. Mã nguồn nền trong dự án sử dụng tám bước. Cách xử lý cố định này đơn giản, nhưng đặt ra câu hỏi: liệu mọi ảnh đều cần dùng đủ tám bước, đặc biệt khi ảnh đầu vào chỉ thiếu sáng ở mức vừa phải?

Từ câu hỏi đó, nhóm xây dựng AdaZero-DCE: giữ nguyên DCE-Net và học một mạng hồi quy nhỏ để dự đoán số bước từ đặc trưng độ sáng. Mục tiêu của bài tập là triển khai một thay đổi có thể kiểm chứng, đo chênh lệch với baseline trên cùng giao thức và phân tích phạm vi mà thay đổi đó có ích. Báo cáo không đặt mục tiêu khẳng định phương pháp tốt hơn mọi mô hình tăng cường ảnh hiện có.

\subsection{Đầu vào, đầu ra và phạm vi}
Đầu vào là ảnh RGB $I\in[0,1]^{H\times W\times3}$. Đầu ra gồm ảnh tăng cường $\widehat I$ có cùng kích thước và số vòng $n\in\{1,\ldots,8\}$. Trong suy luận, hệ thống chỉ quan sát $I$; ảnh tham chiếu đủ sáng $Y$ được sử dụng để đánh giá, không tham gia quyết định số vòng.

Huấn luyện được giới hạn ở head hồi quy. Trọng số DCE-Net \texttt{Epoch99.pth} luôn đóng băng. Dữ liệu có cặp chỉ được dùng để xây dựng nhãn số vòng tối ưu trên tập train. Vì head học từ nhãn được suy ra qua ảnh tham chiếu, toàn bộ phương pháp AdaZero-DCE trong báo cáo là một mở rộng có giám sát, dù backbone xuất phát từ phương pháp zero-reference. Phân biệt này giúp tránh hiểu sai tên gọi của mô hình.

\subsection{Các hướng tiếp cận liên quan}
Các phép biến đổi như tăng gain hoặc gamma có ưu điểm là dễ sử dụng, song mức biến đổi cần được lựa chọn phù hợp. Nhóm phương pháp dựa trên Retinex mô hình hóa quan hệ giữa phản xạ và chiếu sáng; Retinex-Net là một ví dụ học sâu gắn với dữ liệu LOL \cite{lol}. Zero-DCE học bản đồ đường cong mà không yêu cầu cặp ảnh tham chiếu trong cách huấn luyện gốc \cite{zerodce}.

Ý tưởng điều chỉnh thích ứng đã được nghiên cứu trong các công trình khác. VAR-LIDE sử dụng thông tin visibility từ mô hình thị giác-ngôn ngữ để điều kiện hóa quá trình tăng cường đường cong trong bài toán kết hợp tăng sáng và khử mờ \cite{varlide}. AdaZero-DCE trong bài tập này sử dụng đặc trưng thống kê đơn giản và một head nhỏ để chọn tiền tố của tám bước đã có. Nhóm không tuyên bố đây là công trình đầu tiên sử dụng cơ chế thích ứng, và không đưa ra so sánh định lượng với VAR-LIDE vì chưa tái lập công trình đó trong cùng giao thức.

\subsection{Nội dung thực hiện}
Nhóm thực hiện đầy đủ việc tạo 1940 nhãn oracle từ 485 cặp train, huấn luyện ba head với cấu hình cố định và đánh giá fixed-8/adaptive trên cùng implementation. Hai nhóm dữ liệu kiểm tra gồm LOLv1 eval15 với ba điều kiện giả lập bổ sung, và LOLv2-real Test để kiểm tra chuyển miền. Chương trình thực nghiệm đánh giá lại các mô hình đã huấn luyện, lưu kết quả đầu ra, bảng CSV và hình ảnh để truy vết số liệu \cite{artifact}.

\section{Xây dựng và kiểm tra dữ liệu}
\subsection{LOLv1 và phân chia tập dữ liệu}
LOLv1 là dữ liệu ảnh thiếu sáng có ảnh tham chiếu đủ sáng \cite{lol}. Dự án sử dụng nguyên phân chia \texttt{our485} gồm 485 cặp để train head và \texttt{eval15} gồm 15 cặp để đánh giá. Các cặp low/high được ghép theo tên file. Trước khi huấn luyện, pipeline kiểm tra số lượng, tập tên file và mã băm nội dung giữa hai tập để phát hiện ảnh trùng.

Với mỗi cặp train $(I,Y)$, bốn đầu vào được tạo thành:
\begin{equation}
\mathcal{I}(I,Y)=\{I,\;0.3Y,\;0.45Y,\;0.6Y\}.
\end{equation}
Như vậy, số mẫu dùng tính oracle và train head là $485\times4=1940$. Ký hiệu \texttt{factor=1.0} trong artifact là ảnh low gốc $I$, không phải ảnh $Y\times1.0$. Ba hệ số còn lại là ảnh GT nhân trực tiếp trên miền RGB đã chuẩn hóa.

\begin{table}[H]\centering
\caption{Phân chia dữ liệu được sử dụng trong dự án.}
\begin{tabular}{lrrl}\toprule
Tập & Cặp gốc & Đầu vào & Vai trò\\\midrule
LOLv1 our485 & 485 & 1940 & Train head, tính oracle\\
LOLv1 eval15 & 15 & 60 & Đánh giá 4 điều kiện\\
LOLv2-real Test & 100 & 100 & Đánh giá khái quát\\\bottomrule
\end{tabular}
\end{table}

\subsection{Ý nghĩa của các mức sáng giả lập}
Nhân ảnh GT với hệ số nhỏ hơn một tạo ra các trường hợp kiểm soát được về cường độ. Điều này giúp khảo sát việc sử dụng tám vòng có thể vượt quá mức cần thiết khi ảnh chỉ bị giảm cường độ đơn giản. Tuy nhiên, phép nhân trên ảnh sRGB không mô phỏng toàn bộ quá trình chụp ảnh thiếu sáng. Nhiễu cảm biến, sai lệch màu, mất chi tiết do phơi sáng và xử lý trong camera không được tái tạo đầy đủ.

\fig{data_conditions.png}{1}{Một cặp eval15 và các điều kiện giả lập được sử dụng trong thực nghiệm. Ba ảnh giữa được tạo từ cùng GT.}
\textbf{Nhận xét:} Các biến thể của cùng một cặp không phải các cảnh độc lập. Tập train vẫn chỉ có 485 cảnh gốc và tập eval15 vẫn chỉ có 15 cảnh. Việc báo 1940 mẫu không đồng nghĩa tăng số cảnh lên bốn lần. Tương tự, ba seed không làm tăng số ảnh test độc lập.

\subsection{LOLv2-real và kiểm soát nguồn dữ liệu}
LOLv2 là phần mở rộng của LOL, được giới thiệu trong công trình về Deep Retinex Network \cite{lolv2}. Lần đánh giá bổ sung chỉ sử dụng 100 cặp thuộc \texttt{Real\_captured/Test}. Các ảnh Low và Normal được ghép theo mã số, ví dụ \texttt{low00690.png} với \texttt{normal00690.png}. Không tải phần Train, không tạo oracle LOLv2 và không cập nhật tham số từ tập này.

Chương trình chuẩn bị dữ liệu tải riêng 200 ảnh kiểm tra từ bản sao dữ liệu trên Kaggle (\nolinkurl{tanhyml/lol-v2-dataset}), theo danh sách file đã lưu. Tổng dung lượng kiểm tra là 75,970,040 byte. Mỗi ảnh được kiểm tra kích thước byte, khả năng giải mã và lưu SHA256 để truy vết. Đây là kiểm tra tính nhất quán của dữ liệu thực tế đã dùng; chưa đối chiếu checksum của mirror với bản phát hành Google Drive của tác giả. Hạn chế nguồn dữ liệu này được giữ rõ trong báo cáo.

\subsection{Tiền xử lý và tránh rò rỉ dữ liệu}
Ảnh được đọc dưới dạng RGB float32 rồi chia cho 255. Không resize, crop, lượng tử hóa đầu ra trước khi tính metric hoặc điều chỉnh độ sáng đầu ra theo GT. Kích thước low và GT phải khớp. Trong lần chạy này, ảnh giữ độ phân giải gốc $600\times400$.

Chuẩn hóa đặc trưng dùng trung bình và độ lệch chuẩn từ 1940 mẫu train. Cùng các thống kê này được nạp lại khi đánh giá eval15 và LOLv2. Không fit lại normalization trên test. Mọi lựa chọn $n$ được thực hiện trước khi tính PSNR/SSIM với GT. Đây là điều kiện cần để phép đo phản ánh khả năng suy luận thực tế.

\section{Phương pháp AdaZero-DCE}
\subsection{Đường cong tăng cường của Zero-DCE}
Một dạng biểu diễn đường cong thường dùng là $x+\alpha x(1-x)$. Trong implementation đang sử dụng, quy ước dấu tương đương được viết:
\begin{equation}\label{eq:curve}
x_t=x_{t-1}+r_t\odot(x_{t-1}^{2}-x_{t-1}),\qquad x_0=I,
\end{equation}
trong đó $r_t$ là bản đồ hệ số cho bước thứ $t$ và $\odot$ là phép nhân từng phần tử. Hai cách viết tương đương khi $r_t=-\alpha_t$. Do đó, không đổi dấu hệ số chỉ để làm công thức giống cách trình bày trong tài liệu khác.

DCE-Net sinh tám bản đồ RGB từ ảnh ban đầu trong một lần forward. Với $r_t\in[-1,1]$, phép biến đổi duy trì miền giá trị $[0,1]$ về mặt toán học. Đạo hàm theo giá trị điểm ảnh là $1+r_t(2x-1)$, không âm trên miền này khi xem hệ số tại điểm ảnh là cố định. Tính chất đó lý giải vai trò của đường cong trong biến đổi cường độ, nhưng không đảm bảo chất lượng tăng cường tối ưu cho mọi cảnh.

Fixed-8 trả về $x_8$. AdaZero-DCE trả về $x_n$, tức sử dụng $n$ bước đầu tiên của cùng chuỗi bản đồ. Phương pháp không chọn tùy ý một tập con bản đồ, không dùng lại một bản đồ cho mọi bước và không chạy lại DCE-Net tại từng bước.

\subsection{Nhãn oracle trên tập train}
Với mỗi đầu vào train, pipeline tính đủ tám ảnh $x_1,\ldots,x_8$ và chọn:
\begin{equation}
n^*=\underset{n\in\{1,\ldots,8\}}{\operatorname{argmin}}\ \frac{1}{3HW}\|x_n-Y\|_2^2.
\end{equation}
Tối thiểu hóa MSE tương đương tối đa hóa PSNR khi miền giá trị cố định. Implementation tính sai số với float64; nếu nhiều giá trị cực tiểu bằng nhau, thứ tự tìm kiếm chọn chỉ số nhỏ nhất. Oracle chỉ là lựa chọn tốt nhất trong tám ảnh ứng viên do backbone này sinh ra, không phải nghiệm tối ưu cho mọi phương pháp tăng cường ảnh.

Nhãn oracle được tính cho toàn bộ 1940 mẫu. Không áp dụng ngưỡng tối để bỏ qua việc tính nhãn. Quy tắc ngưỡng được sử dụng ở suy luận; cách tổ chức này bảo đảm cache train vẫn phản ánh đầy đủ các ứng viên. Oracle không được tính trên eval15 hoặc LOLv2 trong giao thức đánh giá hiện tại.

\subsection{Đặc trưng độ sáng và head hồi quy}
Tại mỗi điểm ảnh, giá trị dùng cho thống kê là $g=(R+G+B)/3$. Đây là trung bình ba kênh, không phải độ chói có trọng số theo mô hình cảm nhận. Vector đặc trưng gồm:
\begin{equation}
f(I)=[\mu_g,\sigma_g,h_1,\ldots,h_8]\in\mathbb{R}^{10},
\end{equation}
trong đó $h_k$ là tỷ lệ điểm ảnh trong bin thứ $k$ của histogram đều trên $[0,1]$. Tổng tám thành phần histogram bằng một. Chuẩn hóa từng chiều theo tập train:
\begin{equation}
\widetilde f_j=\frac{f_j-\mu_j^{\mathrm{train}}}{\sigma_j^{\mathrm{train}}+10^{-8}}.
\end{equation}
Head gồm ba lớp tuyến tính theo kích thước $10\rightarrow128\rightarrow64\rightarrow1$, ReLU ở hai lớp ẩn và dropout 0.1 sau ReLU thứ nhất. Số tham số là:
\begin{equation}
(10\cdot128+128)+(128\cdot64+64)+(64\cdot1+1)=9729.
\end{equation}
Kích thước này nhỏ so với một mạng xử lý ảnh đầy đủ. Tuy nhiên, đặc trưng toàn cục bỏ qua bố trí không gian: hai ảnh có histogram giống nhau có thể chứa các vùng sáng tối rất khác nhau. Đây là giới hạn biểu diễn quan trọng khi diễn giải kết quả chuyển miền.

\subsection{Hàm mất mát và cấu hình huấn luyện}
Đầu ra liên tục $q_\theta(\widetilde f)$ được hồi quy về nhãn $n^*$. SmoothL1 với tham số mặc định $\beta=1$ có dạng:
\begin{equation}
\ell(e)=\begin{cases}\tfrac12 e^2,&|e|<1,\\|e|-\tfrac12,&|e|\ge1,\end{cases}\qquad e=q_\theta(\widetilde f)-n^*.
\end{equation}
Hàm mất mát trung bình trên toàn bộ tập train được tối ưu bằng Adam. Cấu hình không thay đổi giữa ba seed, được trình bày ở Bảng~\ref{tab:config}.
\begin{table}[H]\centering\caption{Cấu hình cố định của lần train full.}\label{tab:config}
\begin{tabular}{ll}\toprule
Thành phần & Giá trị\\\midrule
DCE-Net & Epoch99, đóng băng\\
Head & 10--128--64--1, 9729 tham số\\
Activation / dropout & ReLU / 0.1\\
Loss & SmoothL1, $\beta=1$\\
Optimizer & Adam, learning rate $5\times10^{-3}$\\
Scheduler & CosineAnnealingLR, $T_{\max}=1200$\\
Epoch / batch & 1200 / full batch 1940\\
Seed & 0, 1, 2\\
Ngưỡng ảnh tối & mean $<0.10$ thì $n=8$\\\bottomrule
\end{tabular}
\end{table}
Head được train trên CPU theo implementation hiện tại; GPU T4 được dùng cho DCE-Net khi tạo oracle và đánh giá. Việc gọi toàn bộ phiên chạy là ``trên Kaggle GPU'' không có nghĩa mọi phép toán của head đều được đặt trên GPU. Checkpoint lưu cả state dict và thống kê chuẩn hóa để suy luận nhất quán.

\subsection{Quy tắc suy luận}
Với một ảnh mới, số vòng được quyết định bởi:
\begin{equation}
n(I)=\begin{cases}
8,&\mu_g<0.10,\\
\operatorname{clip}(\operatorname{round}(q_\theta(\widetilde f)),1,8),&\mu_g\ge0.10.
\end{cases}
\end{equation}
Quy tắc dùng dấu nhỏ hơn nghiêm ngặt; trường hợp mean bằng 0.10 đi qua head. Dropout được tắt bằng chế độ eval. Backbone và head đều không cập nhật gradient trong đánh giá.

\begin{center}
\fbox{\parbox{.88\linewidth}{\centering
Ảnh RGB đầu vào\\$\downarrow$\\
Mean, std, histogram $\longrightarrow$ guard hoặc head $\longrightarrow n$\\[4pt]
Ảnh RGB $\longrightarrow$ DCE-Net đóng băng $\longrightarrow r_1,\ldots,r_8$\\$\downarrow$\\
Áp dụng $n$ bước đầu tiên $\longrightarrow$ ảnh tăng cường
}}
\end{center}

\textbf{Nhận xét:} Guard đảm bảo nhánh ảnh rất tối trả về đúng kết quả fixed-8. Điều này không chứng minh head dự đoán tốt cho các ảnh đó, bởi head đã bị bỏ qua. Khi báo cáo kết quả cần chỉ ra tỷ lệ ảnh đi qua guard để người đọc biết phần nào thực sự kiểm tra head.

\section{Thiết kế thực nghiệm và kết quả}
\subsection{Mục tiêu và giao thức so sánh}
Thực nghiệm kiểm tra ba câu hỏi: điều chỉnh số vòng có giúp các ảnh giả lập thiếu sáng vừa phải hay không; có giữ kết quả trên eval15 gốc hay không; và head LOLv1 có khái quát sang ảnh thật LOLv2 hay không. Zero-DCE baseline (fixed-8) và AdaZero-DCE (adaptive) dùng cùng backbone, ảnh đầu vào, GT và hàm tính metric. Khác biệt duy nhất trong phép so sánh là số bước đường cong được sử dụng.

Không đối chiếu giá trị SSIM tuyệt đối với số kỳ vọng trong kế hoạch hoặc với bảng của công trình dùng implementation khác. Đại lượng chính là chênh lệch ghép cặp giữa hai phương pháp trong cùng lần chạy. Kết quả được tính cho cả ba seed, không chọn seed tốt nhất bằng test.

\subsection{Các chỉ số đánh giá}
Với RGB chuẩn hóa về $[0,1]$, PSNR được tính theo:
\begin{equation}
\operatorname{MSE}(\widehat I,Y)=\frac{1}{3HW}\sum_{p,c}(\widehat I_{p,c}-Y_{p,c})^2,\qquad
\operatorname{PSNR}=-10\log_{10}(\operatorname{MSE}).
\end{equation}
PSNR cao hơn biểu thị sai số bình phương nhỏ hơn đối với cùng GT. Metric được tính cho từng ảnh rồi lấy trung bình; không lấy MSE của toàn bộ tập để đổi sang một PSNR chung.

SSIM đánh giá quan hệ cấu trúc dựa trên thống kê cục bộ \cite{ssim}. Implementation dùng \texttt{skimage.metrics.structural\_similarity}, \texttt{data\_range=1}, \texttt{channel\_axis=2}, cửa sổ đều $7\times7$ và sample covariance \cite{skimage}. Không dùng Gaussian window hoặc chỉ kênh Y. Do lựa chọn cửa sổ, kênh màu và cách chuẩn hóa ảnh có thể thay đổi SSIM, toàn bộ kết luận chỉ dựa trên cặp fixed-8/adaptive được đo nhất quán.

Với metric $M$, chênh lệch là $\Delta M=M_{\mathrm{adaptive}}-M_{\mathrm{fixed8}}$. Giá trị dương biểu thị cải thiện theo metric đó. Trung bình ba seed được tính từ trung bình theo ảnh của từng seed. Độ lệch chuẩn dùng mẫu số 3 (\texttt{ddof=0}); đây là biến thiên theo seed, không phải khoảng tin cậy theo ảnh.

\subsection{Nguồn thực thi của báo cáo}
Quy trình đánh giá gồm các bước chuẩn bị môi trường, kiểm tra dữ liệu, nạp mô hình và tính chỉ số chất lượng. Chương trình sử dụng nguyên ba bộ trọng số từ lần huấn luyện đầy đủ, xác minh mã băm và thực hiện lại quá trình suy luận. Lịch sử hàm mất mát và nhãn oracle được đọc từ kết quả huấn luyện đã lưu, phân biệt với số liệu đánh giá mới. Các bảng kết quả trong báo cáo được tổng hợp từ tệp CSV của lần thực nghiệm này \cite{artifact}.

\fig{training.png}{1}{Hàm mất mát và phân bố nhãn oracle của lần huấn luyện đầy đủ, được dựng lại từ kết quả đã lưu.}
\textbf{Nhận xét:} Phân bố nhãn không đều: số mẫu oracle bằng 8 là 878, trong khi không có mẫu oracle bằng 1. Vì vậy, bài toán không cung cấp lượng thông tin như nhau cho mọi số vòng. Loss train giảm chỉ cho biết khả năng khớp các nhãn đã có; nó không đủ để kết luận chất lượng ảnh test tiếp tục tăng khi train dài hơn.

\subsection{Kết quả trên LOLv1 eval15}
\begin{table}[H]\centering\small
\caption{PSNR (dB) trên eval15: trung bình 3 seed, 15 ảnh/điều kiện.}
\begin{tabular}{lrrrr}\toprule
Điều kiện & Fixed-8 & Adaptive & $\Delta$ & Std seed\\\midrule
Low gốc & 14.8611 & 14.8611 & +0.0000 & 0.0000\\
GT$\times0.3$ & 18.3276 & 18.2442 & -0.0834 & 0.0514\\
GT$\times0.45$ & 17.2052 & 18.2826 & +1.0774 & 0.1615\\
GT$\times0.6$ & 16.4163 & 19.6761 & +3.2598 & 0.0924\\
\bottomrule\end{tabular}
\end{table}
\begin{table}[H]\centering\small
\caption{SSIM trên eval15: trung bình 3 seed, 15 ảnh/điều kiện.}
\begin{tabular}{lrrrr}\toprule
Điều kiện & Fixed-8 & Adaptive & $\Delta$ & Std seed\\\midrule
Low gốc & 0.562653 & 0.562653 & +0.000000 & 0.000000\\
GT$\times0.3$ & 0.891626 & 0.897264 & +0.005638 & 0.000630\\
GT$\times0.45$ & 0.867250 & 0.900064 & +0.032814 & 0.001480\\
GT$\times0.6$ & 0.853672 & 0.914620 & +0.060947 & 0.001649\\
\bottomrule\end{tabular}
\end{table}

\fig{lolv1_deltas.png}{1}{Chênh lệch adaptive so với fixed-8 trên eval15. Thanh sai số thể hiện độ lệch chuẩn qua ba seed.}
\textbf{Nhận xét:} Trên 15 ảnh low gốc, fixed-8 và adaptive đạt cùng PSNR trung bình 14.8611 dB và SSIM 0.562653. Ý nghĩa và giới hạn của kết quả này được phân tích tại Mục~\ref{sec:guard}.

Ở mức GT$\times0.3$, SSIM tăng nhẹ nhưng PSNR giảm. Hai metric phản ánh các khía cạnh khác nhau, nên không thể kết luận adaptive tốt hơn đồng thời theo cả hai chỉ số ở điều kiện này. Ở mức GT$\times0.45$ và GT$\times0.6$, PSNR và SSIM cùng tăng, với mức tăng PSNR lớn hơn ở hệ số 0.6. Đây là bằng chứng phù hợp với giả thuyết giảm số vòng có ích khi biến đổi cường độ đầu vào đơn giản và tương đối gần ảnh đủ sáng.

\fig{comparison_1.0.png}{1}{Ví dụ eval15 gốc: fixed-8 và adaptive tạo cùng đầu ra.}
\fig{comparison_0.3.png}{1}{Ví dụ GT$\times0.3$: cùng một ảnh và một seed để minh họa, không thay cho kết quả trung bình.}
\fig{comparison_0.45.png}{1}{Ví dụ GT$\times0.45$ với đầu vào, fixed-8, adaptive và GT.}
\fig{comparison_0.6.png}{1}{Ví dụ GT$\times0.6$: đối chiếu ảnh đầu ra trong cùng điều kiện.}
Các ví dụ lấy cặp đầu tiên theo thứ tự tên file và dùng seed 0, không chọn ảnh theo mức cải thiện cao nhất. Hình hiển thị được clip về miền hợp lệ để trình bày; bảng metric sử dụng trực tiếp đầu ra float theo giao thức đã nêu. Một ví dụ trực quan chỉ minh họa hành vi, không đại diện cho toàn bộ phân bố test.

\subsection{Kết quả trên LOLv2-real Test}
\begin{table}[H]\centering\small
\caption{PSNR (dB) trên 100 cặp LOLv2-real Test.}
\begin{tabular}{lrrr}\toprule
Seed & Fixed-8 & Adaptive & $\Delta$\\\midrule
0 & 18.0598 & 17.8628 & -0.1969\\
1 & 18.0598 & 17.8759 & -0.1839\\
2 & 18.0598 & 17.8340 & -0.2258\\
Trung bình & 18.0598 & 17.8576 & -0.2022\\
\bottomrule\end{tabular}
\end{table}
Độ lệch chuẩn adaptive qua ba seed: 0.0175.

\begin{table}[H]\centering\small
\caption{SSIM trên 100 cặp LOLv2-real Test.}
\begin{tabular}{lrrr}\toprule
Seed & Fixed-8 & Adaptive & $\Delta$\\\midrule
0 & 0.579800 & 0.581476 & +0.001676\\
1 & 0.579800 & 0.581646 & +0.001846\\
2 & 0.579800 & 0.580941 & +0.001141\\
Trung bình & 0.579800 & 0.581354 & +0.001555\\
\bottomrule\end{tabular}
\end{table}
Độ lệch chuẩn adaptive qua ba seed: 0.000300.


\fig{lolv2_analysis.png}{1}{Phân bố mean RGB, ngưỡng guard và quan hệ giữa độ sáng với chênh lệch PSNR trên LOLv2-real.}
\textbf{Nhận xét:} PSNR trung bình giảm trong cả ba seed, trong khi SSIM tăng rất nhỏ. Có 87/100 ảnh được guard giữ ở $n=8$. Với mỗi seed, 2 ảnh tăng PSNR, 87 ảnh bằng và 11 ảnh giảm. Do phần lớn ảnh không thay đổi, trung bình toàn tập có thể che khuất mức giảm trên nhóm thực sự dùng head.

Trên 13 ảnh không đi qua guard, chênh lệch PSNR trung bình lần lượt xấp xỉ $-1.515$, $-1.414$ và $-1.737$ dB cho seed 0, 1 và 2. Như vậy, kết quả chưa chứng minh head khái quát tốt sang LOLv2-real. Việc SSIM tăng nhẹ không làm mất đi kết quả giảm PSNR, và không đủ để khẳng định cải thiện chất lượng cảm nhận khi chưa có đánh giá tương ứng.

\section{Thảo luận}
\subsection{Ưu điểm của phương pháp}
\textbf{Thích ứng số vòng để hạn chế tăng sáng quá mức.} Với ảnh vừa tối, áp dụng đủ tám vòng có thể làm mức tăng sáng vượt nhu cầu của ảnh. AdaZero-DCE chọn số vòng theo đặc trưng độ sáng, nhờ đó có thể hạn chế hiện tượng này. Trên hai điều kiện giả lập GT$\times0.45$ và GT$\times0.6$, PSNR trung bình ba seed tăng lần lượt khoảng $1.08$ và $3.26$ dB so với Zero-DCE baseline (fixed-8); SSIM cũng tăng ở cả hai điều kiện. Đây là lợi ích định lượng rõ nhất của phương pháp trong thực nghiệm hiện tại.

\textbf{Head nhẹ, không cần VLM hoặc mô hình phụ trợ.} Head chỉ có 9729 tham số, xấp xỉ 10K, và sử dụng mean, std cùng histogram độ sáng. Ngoài DCE-Net nền, phương pháp không cần VLM hay một mạng ngoài để chấm độ sáng. Head có thể huấn luyện trên CPU với chi phí thấp, phù hợp ngân sách vài phút: trong lần train full đã lưu, thời gian thực tế chỉ khoảng $6.49$--$6.76$ giây cho mỗi head với 1200 epoch, khi đã có cache đặc trưng và oracle. Bước tạo oracle là công đoạn riêng và đã dùng GPU T4 trong thí nghiệm này.

\textbf{Bảo toàn kết quả trên LOLv1 gốc.} Trên 15 ảnh eval15 gốc, phương pháp giữ nguyên PSNR và SSIM so với baseline. Phạm vi của ưu điểm này được làm rõ tại Mục~\ref{sec:guard}.

\textbf{Gắn thêm mà không cần train lại baseline.} DCE-Net được đóng băng hoàn toàn; chỉ head hồi quy được học từ oracle-n của tập train. Cơ chế chọn số vòng được bổ sung vào chuỗi đường cong sẵn có, không thay trọng số backbone. Cách tổ chức này thuận tiện cho việc tận dụng checkpoint có sẵn và so sánh trực tiếp với fixed-8.

\subsection{Nhược điểm và phạm vi đóng góp}
\textbf{Mức tăng hiện mới được chứng minh trên ảnh giả lập.} Hai mức tăng $1.08$ và $3.26$ dB được đo trên ảnh GT nhân với hệ số, không phải trên một tập ảnh vừa tối chụp thật. Phép nhân này chưa mô phỏng đầy đủ nhiễu, màu sắc và quá trình thu nhận ảnh. Vì vậy, kết quả chưa chứng minh khả năng tránh tăng sáng quá mức trên ảnh vừa tối thật.

\textbf{Ảnh thật tương đương hoặc hơi kém baseline theo PSNR trung bình toàn tập.} Trên LOLv1 eval15 gốc, phương pháp đạt kết quả bằng fixed-8. Trên LOLv2-real, PSNR trung bình ba seed giảm khoảng $0.20$ dB, trong khi SSIM tăng nhẹ khoảng $0.001555$. Mức giảm toàn tập nhỏ không có nghĩa mọi ảnh chỉ giảm nhẹ: phân tích nhóm 13 ảnh không đi qua guard ở phần kết quả cho thấy một số trường hợp giảm rõ hơn.

\textbf{Chưa khái quát tốt ra ngoài phân bố LOLv1.} Head và thống kê chuẩn hóa hiện được học từ LOLv1, gồm ảnh low gốc cùng các mức sáng giả lập. Kết quả trên LOLv2-real cho thấy hiệu quả chuyển sang tập dữ liệu khác còn hạn chế. Chưa có bằng chứng thực nghiệm để khẳng định head hoạt động tốt trên các phân bố ảnh khác.

\textbf{Ý tưởng thích ứng đã có trong công trình trước.} Cơ chế điều chỉnh đường cong thích ứng đã xuất hiện trong VAR-LIDE \cite{varlide}. Do đó, đóng góp của bài tập là một cách triển khai nhẹ dựa trên thống kê độ sáng, head có giám sát và bảng đo định lượng độc lập so với fixed-8 trong cùng implementation. ``Độc lập'' ở đây nghĩa là tự thực hiện phép đo kiểm chứng của bài tập, không phải tuyên bố ý tưởng thích ứng hoàn toàn mới hoặc đã tốt hơn VAR-LIDE.

\subsection{Khác biệt hiệu quả giữa ảnh giả lập và ảnh thật}
Ba điều kiện giả lập được tạo bằng nhân trực tiếp GT với một hệ số. Chúng giữ nguyên cấu trúc cảnh và nhiều đặc trưng của GT; do đó, vấn đề điều chỉnh cường độ tương đối đơn giản hơn ảnh chụp thiếu sáng thật. Trong trường hợp này, sử dụng ít bước hơn có thể giúp tránh điều chỉnh vượt mức cần thiết. Kết quả tăng đồng thời PSNR/SSIM ở hệ số 0.45 và 0.6 phù hợp với cách giải thích đó.

Ảnh thật có thể khác về nhiễu, màu sắc, chiếu sáng không đồng đều và xử lý của camera. Vector 10 chiều không mã hóa những yếu tố này một cách trực tiếp. Hai ảnh có mean gần nhau không nhất thiết cần cùng số bước. Đây là một giải thích khả dĩ cho kết quả LOLv2, nhưng chưa phải kết luận nhân quả vì dự án chưa có ablation riêng cho từng yếu tố.

\subsection{Vai trò và giới hạn của guard}\label{sec:guard}
Ngưỡng 0.10 giúp tránh thay đổi những ảnh rất tối theo thống kê toàn cục. Trên eval15 gốc, nó phủ toàn bộ tập; trên LOLv2, nó phủ 87\%. Do đó, mức ``không giảm'' trên eval15 chỉ chứng minh quy tắc đã bảo toàn baseline ở những ảnh đó. Không nên sử dụng kết quả này để khẳng định head đúng cho mọi ảnh thiếu sáng.

Ngưỡng được giữ nguyên từ kế hoạch. Nhóm không chọn lại ngưỡng bằng LOLv2 Test sau khi quan sát kết quả xấu. Nếu cần nghiên cứu ngưỡng tốt hơn, phải xây dựng tập validation riêng trong train và khóa cấu hình trước lần đánh giá cuối. Việc chỉnh trực tiếp theo test sẽ làm phép đo sau đó mất tính độc lập.

\subsection{Ảnh hưởng của thời gian huấn luyện đến chất lượng ảnh}
Kéo dài huấn luyện không bảo đảm PSNR/SSIM trên tập kiểm tra tăng. SmoothL1 tối ưu dự đoán nhãn oracle, còn kết quả cuối phụ thuộc vào bước làm tròn, guard và mức sai khác giữa các ảnh ứng viên. Một thay đổi nhỏ ở đầu ra liên tục có thể không đổi $n$, hoặc làm $n$ nhảy sang số nguyên khác. Vì thế, loss giảm không dẫn đến PSNR/SSIM tăng đơn điệu.

Lần train full đã tuân thủ 1200 epoch cho 1940 mẫu và ba seed. Kết quả hiện tại cho thấy hiệu quả phụ thuộc miền dữ liệu, thay vì chỉ phụ thuộc số epoch. Chưa có thí nghiệm sweep số epoch với validation độc lập, nên báo cáo không suy diễn rằng kéo dài huấn luyện sẽ giải quyết hạn chế chuyển miền. Tương tự, ba seed ổn định về dấu chênh lệch không thay thế cho đánh giá trên nhiều bộ dữ liệu.

\subsection{Chi phí tính toán và tính công bằng}
Số bước đường cong có thể giảm, nhưng DCE-Net vẫn sinh đủ tám bản đồ trong một forward. Vì vậy, không thể đồng nhất giảm từ tám xuống bốn bước với tăng tốc toàn hệ thống gấp đôi. Phần tiết kiệm chỉ nằm ở các phép áp dụng đường cong; head và trích đặc trưng còn bổ sung chi phí. Chưa có benchmark latency đồng bộ GPU trên phần cứng triển khai, nên báo cáo không đưa ra tỷ lệ tăng tốc.

Quy trình đánh giá dùng chung các ảnh ứng viên để tính chỉ số cho cả fixed-8 và adaptive, giúp so sánh chất lượng trên cùng backbone. Vì vậy, thời gian thực hiện toàn bộ quy trình không đại diện cho thời gian suy luận của một lần tăng cường ảnh trong ứng dụng. Muốn đánh giá hiệu năng cần đo riêng đường chạy triển khai, warm-up, đồng bộ thiết bị và thống kê nhiều lần chạy.

\subsection{Giới hạn đánh giá bổ sung}
Tập eval15 nhỏ và ba ảnh giả lập của một cảnh có tương quan mạnh. Độ lệch chuẩn qua ba seed chỉ đo độ nhạy đối với khởi tạo head; chưa có bootstrap theo cảnh hoặc kiểm định thống kê. Vì vậy, các chênh lệch nhỏ cần được diễn giải thận trọng.

Nhãn oracle tối ưu PSNR, không tối ưu trực tiếp SSIM hoặc độ dễ quan sát đối với con người. Head dùng đặc trưng toàn cục và lựa chọn một số vòng chung cho toàn ảnh. DCE-Net được giữ nguyên nên khả năng sửa nhiễu, màu và chi tiết vẫn bị giới hạn bởi backbone. Ngoài ra, nguồn LOLv2 mirror chưa được đối chiếu checksum với bản tác giả. Các hạn chế này xác định phạm vi của kết luận, không được che bởi kết quả thuận lợi trên dữ liệu giả lập.

\section{Tổ chức chương trình và khả năng sử dụng}
\subsection{Các thành phần trong repository}
\begin{table}[H]\centering\small
\caption{Vai trò các thành phần hiện có và artifact bổ sung.}
\begin{tabularx}{\linewidth}{lX}\toprule
Thành phần & Vai trò\\\midrule
\texttt{zero\_dce\_adaptive\_n.py} & Đặc trưng, head, oracle, suy luận thích ứng.\\
\texttt{experiment.py} & Train head và đánh giá LOLv1.\\
\texttt{evaluate\_lolv2.py} & Kiểm tra test LOLv2 bằng checkpoint đóng băng.\\
\texttt{infer.py} & Tăng cường một ảnh bằng head đã lưu.\\
\texttt{AdaZero-DCE\_code/} & Backbone và tài nguyên nền kế thừa Zero-DCE.\\
\texttt{notebook/} & Chương trình thực nghiệm, kết quả thực thi và bảng CSV.\\
\texttt{report/} & Nguồn LaTeX, bảng và hình của báo cáo.\\\bottomrule
\end{tabularx}
\end{table}
Trong giai đoạn hoàn thiện báo cáo, mã nguồn hiện có và checkpoint được giữ nguyên. Các tên Zero-DCE thuộc bài báo, checkpoint, mã nền hoặc tài nguyên tham chiếu vẫn giữ đúng nguồn gốc. Tên AdaZero-DCE dùng cho phương pháp mở rộng và tài liệu hiện hành; không thay tên công trình gốc để tạo cảm giác toàn bộ repository do nhóm xây dựng mới.

\subsection{Tổ chức chương trình thực nghiệm và luồng sử dụng}
Chương trình thực nghiệm được tổ chức dưới dạng tệp Jupyter Notebook, gồm mã nguồn và kết quả thực thi đi kèm. Nội dung lần lượt trình bày môi trường, dữ liệu, mô hình, lịch sử huấn luyện, đánh giá trên LOLv1 và LOLv2, cùng các bước kiểm tra tính tái lập. Các khối lệnh sử dụng chung mã triển khai và hàm tính chỉ số của dự án. Bảng kết quả được xuất ra CSV để có thể kiểm tra độc lập với phần diễn giải.

Đối với một ảnh sử dụng thực tế, chương trình đọc RGB, kiểm tra miền pixel, nạp backbone cùng head và thống kê train, dự đoán $n$ rồi áp dụng các bước đường cong. Không cần GT. Mô hình này có thể dùng làm thành phần thử nghiệm tăng sáng, nhưng kết quả LOLv2 cho thấy chưa đủ cơ sở để khuyến nghị thay fixed-8 cho mọi ảnh thật.

\section{Kết luận và hướng phát triển}
Bài tập đã triển khai cơ chế điều chỉnh số vòng trên backbone Zero-DCE đóng băng, huấn luyện đầy đủ 485 cặp với bốn điều kiện thành 1940 mẫu, giữ nguyên cấu hình 1200 epoch và chạy ba seed. Quy trình đánh giá lại sử dụng các trọng số đã lưu, cung cấp kết quả thực thi, bảng CSV và hình minh họa có thể truy vết.

Điểm mạnh của phương pháp là lựa chọn số vòng phù hợp với độ sáng để hạn chế tăng sáng quá mức trên ảnh vừa tối, sử dụng head khoảng 10K tham số và không cần train lại DCE-Net. Trên hai điều kiện giả lập GT nhân 0.45 và 0.6, PSNR tăng lần lượt khoảng 1.08 và 3.26 dB so với fixed-8. Ngưỡng bảo vệ giữ kết quả không tụt trên LOLv1 eval15 gốc.

Giới hạn chính là mức tăng mới được đo trên ảnh giả lập; trên ảnh thật, kết quả bằng baseline ở LOLv1 và hơi kém theo PSNR trung bình toàn tập ở LOLv2-real (giảm khoảng 0.20 dB, dù SSIM tăng nhẹ). Head hiện học theo phân bố LOLv1 và chưa khái quát tốt sang LOLv2-real. Đóng góp của bài tập nằm ở cách làm nhẹ và bảng định lượng tự thực hiện; ý tưởng thích ứng cần được đặt trong mối liên hệ với VAR-LIDE, không tuyên bố là ý tưởng lần đầu xuất hiện.

Hướng tiếp theo là thiết kế validation độc lập, kiểm tra đặc trưng có thông tin không gian hoặc nhiễu, nghiên cứu mục tiêu lựa chọn số vòng phù hợp hơn và đo latency triển khai. Đây là các đề xuất nghiên cứu tiếp, chưa được thực hiện trong báo cáo. Mọi thay đổi cần được quyết định bằng train/validation và giữ tập test cho đánh giá cuối, thay vì chỉnh mô hình theo các kết quả test đã quan sát.

\clearpage
\begin{thebibliography}{9}
\bibitem{zerodce} C. Guo và cộng sự, ``Zero-Reference Deep Curve Estimation for Low-Light Image Enhancement,'' CVPR, 2020. \url{https://github.com/Li-Chongyi/Zero-DCE}.
\bibitem{lol} C. Wei, W. Wang, W. Yang và J. Liu, ``Deep Retinex Decomposition for Low-Light Enhancement,'' BMVC, 2018. \url{https://daooshee.github.io/BMVC2018website/}.
\bibitem{lolv2} W. Yang, H. Huang, W. Wang, S. Wang và J. Liu, ``Sparse Gradient Regularized Deep Retinex Network for Robust Low-Light Image Enhancement,'' IEEE Transactions on Image Processing, 2021. \url{https://github.com/flyywh/CVPR-2020-Semi-Low-Light}.
\bibitem{varlide} W. Dong, H. Zhou, J. Lin và J. Chen, ``Zero-Reference Joint Low-Light Enhancement and Deblurring via Visual Autoregressive Modeling with VLM-Derived Modulation,'' arXiv:2511.18591, 2025. \url{https://arxiv.org/abs/2511.18591}.
\bibitem{ssim} Z. Wang, A. C. Bovik, H. R. Sheikh và E. P. Simoncelli, ``Image Quality Assessment: From Error Visibility to Structural Similarity,'' IEEE Transactions on Image Processing, 13(4), 600--612, 2004.
\bibitem{skimage} scikit-image, tài liệu \texttt{skimage.metrics.structural\_similarity}. \url{https://scikit-image.org/docs/stable/api/skimage.metrics.html}. Truy cập ngày 19/10/2026.
\bibitem{artifact} Nhóm thực hiện, ``AdaZero DCE Complete Evaluation Report,'' chương trình thực nghiệm trên Kaggle và các bảng CSV đi kèm, ngày 19/10/2026. \url{https://www.kaggle.com/code/sangtranminh211/adazero-dce-complete-evaluation-report}. Chương trình được lưu ở chế độ riêng tư; bản kèm kết quả thực thi nằm trong thư mục \texttt{notebook/} của dự án.
\end{thebibliography}
\clearpage
\appendix
\section{Bảng kết quả chi tiết và khả năng tái lập}
\begin{table}[H]\centering\small
\caption{Kết quả chi tiết từng seed trên eval15. Baseline xem bảng chính.}
\begin{tabular}{lrrrrrr}\toprule
Điều kiện & Seed & PSNR & $\Delta$PSNR & SSIM & $\Delta$SSIM & $\bar n$\\\midrule
Low gốc & 0 & 14.8611 & +0.0000 & 0.562653 & +0.000000 & 8.000\\
Low gốc & 1 & 14.8611 & +0.0000 & 0.562653 & +0.000000 & 8.000\\
Low gốc & 2 & 14.8611 & +0.0000 & 0.562653 & +0.000000 & 8.000\\
GT$\times0.3$ & 0 & 18.3133 & -0.0143 & 0.897251 & +0.005625 & 7.400\\
GT$\times0.3$ & 1 & 18.1901 & -0.1374 & 0.898042 & +0.006416 & 7.267\\
GT$\times0.3$ & 2 & 18.2291 & -0.0984 & 0.896499 & +0.004873 & 7.333\\
GT$\times0.45$ & 0 & 18.1534 & +0.9482 & 0.899351 & +0.032101 & 6.600\\
GT$\times0.45$ & 1 & 18.5104 & +1.3052 & 0.902125 & +0.034875 & 6.733\\
GT$\times0.45$ & 2 & 18.1841 & +0.9789 & 0.898717 & +0.031467 & 6.667\\
GT$\times0.6$ & 0 & 19.7406 & +3.3244 & 0.914977 & +0.061305 & 6.133\\
GT$\times0.6$ & 1 & 19.7421 & +3.3259 & 0.916436 & +0.062764 & 6.067\\
GT$\times0.6$ & 2 & 19.5455 & +3.1292 & 0.912446 & +0.058773 & 6.133\\
\bottomrule\end{tabular}
\end{table}

\subsection{Quy trình chạy lại}
Cài thư viện theo \texttt{requirements.txt}. Để tái lập huấn luyện head, dùng \texttt{experiment.py} với \texttt{--train-pairs 485 --epochs 1200 --seeds 0 1 2}. Đây là lệnh huấn luyện; quy trình đánh giá lại không thực hiện bước này. Để kiểm tra kết quả, sử dụng chương trình thực nghiệm được cung cấp cùng hai nguồn kết quả trên Kaggle của lần huấn luyện đầy đủ và lần đánh giá LOLv2 trước đó.

Quy trình tái lập lần lượt kiểm tra SHA256, tải dữ liệu kiểm tra, chạy suy luận, tính các chỉ số chất lượng, xuất CSV và đối chiếu với kết quả đã lưu trong sai số tuyệt đối $10^{-6}$. Mọi khối lệnh phải thực thi thành công; khối kiểm tra cuối cùng phải trả về thông báo PASS. Tệp kèm kết quả đầu ra được tạo bằng \texttt{nbclient} trên Kaggle rồi tải về, không chèn số liệu thủ công.

Các bảng trong báo cáo được sinh từ CSV của lần chạy này. Manifest lưu seed, mã băm checkpoint, phiên bản PyTorch, GPU và số bản ghi. Lịch sử loss được đọc từ lần train full trước đó; không diễn giải đồ thị này là một lần huấn luyện mới. Thời gian thực hiện quy trình đánh giá bao gồm suy luận, xử lý dữ liệu và tính chỉ số chất lượng; số đo này không đại diện cho độ trễ triển khai.
\end{document}
